\documentclass[a4paper,11pt]{article}
\usepackage[misc]{ifsym}
\usepackage{xcolor}

\usepackage[activate={true,nocompatibility}]{microtype}
\usepackage{microtype}
\usepackage[T1]{fontenc}
\usepackage{amsthm}
\newtheorem{lemma}{Lemma}[section]
%\usepackage{mlmodern}
%\usepackage{Baskervaldx}
\usepackage[xcharter]{newtx}
\usepackage{fontawesome5}


\renewcommand{\baselinestretch}{1.06}
\newcommand{\subt}{ <: }
\newcommand{\kmod}{\normalfont\textsf{\textbf{mod}}}
%\usepackage{MnSymbol}
\usepackage{inconsolata}
\usepackage[scaled=0.92]{sourcesanspro}
\usepackage{xcolor}
\usepackage[utf8]{inputenc}
\usepackage[english]{babel}
\usepackage[a4paper, left=1.33in, right=1.33in]{geometry}
\usepackage{calc}
\usepackage{graphicx}

\usepackage{titlesec}
\titleformat{\subsubsection}[runin]
  {\normalfont\itshape}
  {\hspace*{\parindent}\normalfont\bfseries(\thesubsubsection)}
  {1.5\wordsep}
  {}
  []

\titlespacing*{\section}{0pt}{2\baselineskip}{1\baselineskip}
\titlespacing*{\subsubsection}
  {0pt}{6pt}{\wordsep}  % left, before, after spacing

\usepackage{mathpartir}
\usepackage{amsmath}

\usepackage{tikz}
\usepackage{wrapfig}

\usepackage[backend=bibtex, style=alphabetic]{biblatex}
\addbibresource{notes.bib}
\renewcommand*{\bibfont}{\small}

\author{Aghilas Y. Boussaa \texttt{<aghilas.boussaa@ens.fr>}}

\title{Notes on modal effect types}
\begin{document}
\sloppy
\maketitle

\section{Adapting \emph{Complete and Easy Bidirectional Typechecking for Higher-Rank Polymorphism}}
\subsection{Introduction}
\subsubsection{Brief overwiew of the plan.}
We remove explicit type application from \textsf{METL}. We want to adapt the
algorithm of Dunfield and Krishnaswami~\cite{dunfield2013complete} that allows
instantiation with monotypes. I will start by considering \emph{intuitionnistic}
monotypes (i.e.\ without modalities) to start, guarded monotypes may actually
work.

\subsubsection{Why monotypes?}
The paper mentions in Section~2 that it sticks to predicative polymorphism for
monotypes as the subtyping relation for general system F is undecidable.
However, it may still be decidable by allowing modalities. Allowing modalities
also means that we need to insert boxing and unboxing, when substituting a type
with one that has a top-level modality.

\subsubsection{Which subtyping relation?}
The original algorithm relies heavily on a subtyping relation, which we shall
adapt to \textsf{MET} types. To support unification from \textsf{METL} to
\textsf{MET}, the subtyping algorithm also has an elaboration part which
rewrites gives a function from $A$ to $B$ when $A$ is a subtype of $B$:
$\Gamma\vdash A\textsf{ <: }B\dashv\Delta{\color{gray} \dashrightarrow M}$.

Support for modalities in subtyping will come later. I start by only allowing
different modalities on the top-level. And using equality of modalities inside
types when unifying.

\subsection{Copying the original subtyping relation}
\subsubsection{Contexts.}
We merge the definitions of contexts given by both papers in
Figure~\ref{contexts}. We use the hole notation, meaning that: $\Gamma[\Delta]$
is a context of the form: $\Gamma_L,\Delta,\Gamma_R$.
\begin{figure}[h]
  \[
  \Gamma,\Delta\;::=\;\cdot\mid \Gamma,\alpha:K \mid \Gamma,\hat{\alpha}:K \mid
  \Gamma,\hat{\alpha}=\tau \mid \Gamma, x : A
  \mid \Gamma,\textsf{\scriptsize\faLock}_\mu
  \mid \Gamma, \blacktriangleright_{\hat{\alpha}}
  \]
  \caption{Syntax of contexts}\label{contexts}
\end{figure}


\subsubsection{Subtyping and instantiation.}
We have the same algorithmic subtyping and instantiation rules to which we add
the rule \textsc{Sub-Modality} to cross modalities when they are the same on
both sides; there is no rule for instantiation with modalities. And each rule
has now an elaboration part in which we may omit type annotations when they are
obvious.

When writing rules, we omit the kind of type variables when they are not
significant. In absence of kind annotation, a same variable should be considered
of having the same kind through the whole rule.

\begin{figure}[h]
  \centering
  \begin{mathpar}
    \inferrule[Sub-Var]{\\}{\Gamma[\alpha]\vdash\alpha\subt\alpha\dashv\Gamma[\alpha]
      {\color{gray} \dashrightarrow \lambda x.x}}
    
    \inferrule[Sub-ExVar]{\\}{\Gamma[\hat{\alpha}]\vdash\hat{\alpha}\subt\hat{\alpha}\dashv\Gamma[\hat{\alpha}]
      {\color{gray} \dashrightarrow \lambda x.x}}

    \inferrule[Sub-Unit]{\\}{\Gamma\vdash 1\subt1\dashv\Gamma
      {\color{gray} \dashrightarrow \lambda x.x}}

    \inferrule[Sub-Fun]{\Gamma\vdash B_1\subt A_1\dashv\Theta
      {\color{gray} \dashrightarrow M}\\
      \Theta{\vdash} [\Theta]A_2 {\subt} [\Theta]B_2\dashv\Delta
      {\color{gray} \dashrightarrow N}\\
      M' = \lambda f^{A_1\to A_2} x^{B_1}.N (f(M x))}{
      \Gamma\vdash A_1\to A_2\subt B_1\to B_2\dashv \Delta {\color{gray} \dashrightarrow M'}}

    \inferrule[Sub-ForallL]{\Gamma,\blacktriangleright_{\hat{\alpha}},{\hat{\alpha}:K}
      {\vdash}[{\hat{\alpha}}/\alpha]A\subt B\dashv\Delta,\blacktriangleright_{\hat{\alpha}},\Theta
      {\color{gray} \dashrightarrow M}\\
      {\color{gray} \tau = [\Theta[\hat{\beta}:K/\hat{\beta}=1]]\hat{\alpha}}
      }
      {\Gamma\vdash\forall\alpha:K.A\subt B\dashv\Delta
        {\color{gray} \dashrightarrow \lambda x^{\forall\alpha:K.A}.M(x \tau)}}

  \inferrule[Sub-ForallR]{\Gamma,\alpha:K\vdash A\subt B\dashv\Delta,\alpha:K,\Theta
    {\color{gray} \dashrightarrow M}}
    {\Gamma\vdash A\subt \forall\alpha:K.B\dashv\Delta
      {\color{gray} \dashrightarrow \lambda^A x.\Lambda \alpha^K. M x}}

  \inferrule[Sub-InstantiateL]{\hat{\alpha}\not\in\textsf{FV}(A)\\
    {\Gamma}[\hat{\alpha}]\vdash\hat{\alpha}:\leqq A\dashv\Delta
    {\color{gray} \dashrightarrow M}}
    {{\Gamma}[\hat{\alpha}]\vdash\hat{\alpha}\subt A\dashv\Delta
      {\color{gray} \dashrightarrow M}}

  \inferrule[Sub-InstantiateR]{\hat{\alpha}\not\in\textsf{FV}(A)\\
    {\Gamma}[\hat{\alpha}]\vdash A\leqq:\hat{\alpha}\dashv\Delta
    {\color{gray} \dashrightarrow M}}
    {{\Gamma}[\hat{\alpha}]\vdash A\subt \hat{\alpha}\dashv\Delta
      {\color{gray} \dashrightarrow M}}

  \inferrule[Sub-Modality]{\Gamma\vdash A\subt B\dashv\Delta{\color{gray} \dashrightarrow M}}
    {\Gamma\vdash\mu A\subt\mu B\dashv\Delta{\color{gray} \dashrightarrow \lambda x^A . (\lambda y^B. \kmod_\mu y) (M x) }}
  \end{mathpar}

  \caption{Algorithmic subtyping: $\Gamma\vdash A\subt B\dashv\Delta{\color{gray} \dashrightarrow M}$}
\end{figure}

The adaptation is straightforward expect for rule \textsc{Sub-ForallL} and its
type $\tau = [\Theta[\hat{\alpha}:K/\hat{\alpha}=1]]\hat{\alpha}$. The
substitution replaces all trailing existentials with unit types as their
definition is not important. For instance, when cheking
$\forall\beta:\textsf{Any}.1<:1$, the rule inserts an existential variable whose
actual definition is useless. The type $\tau$ is an arbitrary solution.
\begin{figure}[h]
  \begin{mathpar}
    \inferrule[InstL-Solve]{\Gamma\vdash\tau:K}
      {\Gamma,\hat{\alpha}:K,\Gamma'\vdash\hat{\alpha}:\leqq\tau\dashv\Gamma,\hat{\alpha}=\tau,\Gamma'
      {\color{gray} \dashrightarrow \lambda x.x }}

    \inferrule[InstL-Reach]{\\}
      {{\Gamma}[\hat{\alpha}:K][\hat{\beta}:K'] \vdash\hat{\alpha}:\leqq\hat{\beta}\dashv\Gamma[\hat{\alpha}:K\wedge K'][\hat{\beta}=\hat{\alpha}]
        {\color{gray} \dashrightarrow\lambda x.x }}

    \inferrule[InstL-Arr]
      {{\Gamma}[{\hat{\alpha}}_2 :\textsf{Any}, {\hat{\alpha}}_1 :\textsf{Any},
                \hat{\alpha}={\hat{\alpha}}_1\to{\hat{\alpha}}_2]
        \vdash A_1\leqq: {\hat{\alpha}}_1 \dashv\Theta
        {\color{gray}\dashrightarrow M}\\
        \Theta\vdash {\hat{\alpha}}_2:{\leqq}[\Theta]A_2\dashv\Delta
        {\color{gray}\dashrightarrow N}}
      {{\Gamma}[{\hat{\alpha}:\textsf{Any}}]\vdash\hat{\alpha}:\leqq A_1 \to A_2\dashv\Delta
        {\color{gray} \dashrightarrow \lambda f^{{\hat{\alpha}}_1{\to} {\hat{\alpha}}_2} x^{A_1}.N (f (M x))}}

    \inferrule[InstL-AllR]
      {{\Gamma}[\hat{\alpha}],\beta:K\vdash\hat{\alpha}:\leqq B\dashv\Delta,\beta:K,\Delta'
        {\color{gray} \dashrightarrow M }}
      {{\Gamma}[\hat{\alpha}]\vdash\hat{\alpha}:\leqq \forall\beta:K.B\dashv\Delta
        {\color{gray} \dashrightarrow \lambda x^{\hat{\alpha}}.\Lambda\beta^{K}.M x}}
  \end{mathpar}

\begin{mathpar}
    \inferrule[InstR-Solve]{\Gamma\vdash\tau:K}
      {\Gamma,\hat{\alpha}:K,\Gamma'\vdash\tau\leqq:\hat{\alpha}\dashv\Gamma,\hat{\alpha}=\tau,\Gamma'
      {\color{gray} \dashrightarrow \lambda x.x }}

    \inferrule[InstR-Reach]{\\}
      {{\Gamma}[\hat{\alpha}:K][\hat{\beta}:K'] \vdash\hat{\beta}\leqq:\hat{\alpha}\dashv\Gamma[\hat{\alpha}:K\wedge K'][\hat{\beta}=\hat{\alpha}]
        {\color{gray} \dashrightarrow\lambda x.x }}

    \inferrule[InstR-Arr]
      {{\Gamma}[{\hat{\alpha}}_2:\textsf{Any}, {\hat{\alpha}}_1 : \textsf{Any},
                 \hat{\alpha}={\hat{\alpha}}_1\to{\hat{\alpha}}_2]
        \vdash {\hat{\alpha}}_1:\leqq A_1 \dashv\Theta
        {\color{gray}\dashrightarrow M}\\
        \Theta\vdash [\Theta]A_2{\leqq}:{\hat{\alpha}}_2\dashv\Delta
        {\color{gray}\dashrightarrow N}}
      {{\Gamma}[{\hat{\alpha}}:\textsf{Any}]\vdash A_1 \to A_2\leqq:\hat{\alpha}\dashv\Delta
        {\color{gray} \dashrightarrow \lambda f^{A_1\to A_2} x^{{\hat{\alpha}}_1}.N (f (M x))}}

    \inferrule[InstR-AllL]
      {{\Gamma}[{\hat{\alpha}}],\blacktriangleright_{\hat{\beta}},{\hat{\beta}:K}
      {\vdash}[{\hat{\beta}}/\beta]B\leqq: \hat{\alpha}\dashv\Delta,\blacktriangleright_{\hat{\beta}},\Delta'
      {\color{gray} \dashrightarrow M}\\
      {\color{gray} \tau = [[\hat{\gamma}:K/\hat{\gamma}=1]\Delta']\hat{\beta}}
      }
      {{\Gamma}[\hat{\alpha}]\vdash\forall\beta:K.B\leqq:\hat{\alpha}\dashv\Delta
        {\color{gray} \dashrightarrow \lambda x^{\forall\beta:K.B}. M (x \tau) }
        }
  \end{mathpar}

  \caption{Instantiation}
\end{figure}

The following lemma guarantees the soundness of elaboration.
\begin{lemma}
  If $\Gamma\vdash A\subt B\dashv\Delta{\color{gray} \dashrightarrow f}$, then
  $\Delta\vdash \kmod_{[]} [\Delta]M : []([\Delta]A\to [\Delta]B) \, @ \, E$.
\end{lemma}

\subsubsection{New inference rules.}
We take the rules of \textsf{METL} (Appendix~D of~\cite{tang2025modal}) as a
starting point, remove the rules of both application and type application and
try to add the new inference rule for application --- I will try to add
inference for abstraction later.
\begin{figure}
  \begin{mathpar}

  \end{mathpar}

  \caption{Algorithmic typing}
\end{figure}

\subsection{Subtyping with modalities}
\subsubsection{Problems.}
Allowing modalities in subtyping requires the subtyping relationship to depend
on an environment context: $\Gamma\vdash A\subt B \dashv \Delta \, @ \,E
{\color{gray} \dashrightarrow f}$. When unifying a flexible type variable
$\hat{\alpha}$ against a non-intuitionistic monotype $\tau$, producing the
equation $\hat{\alpha}=\tau$ is not complete. For instance, we may check
$\hat{\alpha}$ first against the type $[\textsf{Gen int}](1\to1)$, and then
against $<\textsf{Gen int}>(1\to1)$.

A complete algorithm may not exist, but we may do better than plain assignation
by recording inequalities modulo effect context in the context:
$\hat{\alpha}\leq \tau\,@\,E$. Elaboration becomes more involved as there are
now holes with functions that can only be filled when a definition for the type
variable is chosen.

\end{document}
